\documentclass[../principal.tex]{subfiles}
\graphicspath{{\subfix{../imagenes/}}}

\begin{document}
\chapter{Potencia}

\section{Definición}

La potencia eléctrica se define como la cantidad de energía eléctrica que se consume o se produce en un circuito eléctrico, por unidad de tiempo.

\section{Unidad de medida}

La unidad de medida de energía eléctrica es el joule (\unit{\joule}), y la unidad de medida de tiempo es el segundo (\unit{\second}).

La unidad de medida de potencia eléctrica entonces es:

\begin{equation}
    \unit{\joule\per\second}
\end{equation}

Por conveniencia se usa el watt (\unit{\watt}), donde:

\begin{equation}
    \unit{\watt} = \unit{\joule\per\second}
\end{equation}


\end{document}


